\section*{Newton polygon at $d=4$ and the universal $\xi_0$ identity (Phase Q1-3)}

\subsection*{Operator}

For a quartic PCF with $a_n=1$ and $b(n)=a_4 n^4 + a_3 n^3 + a_2 n^2
+a_1 n + a_0$, the partial-denominator generating function
$f(z)=\sum_{n\ge 0} Q_n z^n$ satisfies the inhomogeneous linear ODE
$L\,f = (\text{polynomial in } z\text{ of degree }\le 1)$ with
\[
L\ =\ 1\ -\ z\,B(\theta+1)\ -\ z^2,\qquad
B(t)\ =\ a_4 t^4 + a_3 t^3 + a_2 t^2 + a_1 t + a_0,\qquad
\theta\ =\ z\,\partial_z.
\]
The same construction at degree $d$ gives
$L_d=1-z\,B_d(\theta+1)-z^2$.

\subsection*{Newton polygon at $z=0$}

Nonzero monomial points $(i,j)\in\mathbb{Z}_{\ge 0}^2$ of $L$:
$(0,0)$ with coefficient $+1$, $(1,j)$ for $j\in\{0,1,2,3,4\}$ with
coefficients drawn from $-B(t+1)$, and $(2,0)$ with coefficient $-1$.

The lower-left convex hull at $z=0$ has the single nontrivial edge
\[
E:\ (0,0)\ \longrightarrow\ (1,4)
\qquad\text{slope } 1/d = 1/4,
\]
so $L$ has a single irregular singularity at $z=0$ of Gevrey class $d=4$
in $z$, equivalently Gevrey-$1$ in the variable $u$ defined by $z=u^d$.

\subsection*{Characteristic equation along $E$}

Substituting the WKB ansatz $f\sim\exp(c/u)$ with $z=u^4$ and computing
$\theta\,\exp(c/u) = -(c/(d u))\,\exp(c/u)$ gives
\[
\chi(c)\ =\ 1\ -\ \frac{a_4}{4^4}\,c^4\ =\ 1\ -\ \frac{a_4}{256}\,c^4,
\]
with positive real root
\[
\boxed{\ \xi_0(b)\ =\ \frac{4}{a_4^{1/4}}\ =\ \frac{d}{\beta_d^{1/d}}\Big|_{d=4}.\ }
\]

\subsection*{Empirical verification (5+ representatives)}

\begin{tabular}{lrll}
\hline
representative & $a_4$ & $\xi_0$ (50 dp) & error vs.\ $4/a_4^{1/4}$ \\
\hline
mixed\_$S_4$ smallest $|\Delta_4|$ & $1$ & $4.000000\ldots$ & $0$ \\
mixed\_$S_4$ largest  $|\Delta_4|$ & $1$ & $4.000000\ldots$ & $0$ \\
mixed\_$D_4$                       & $1$ & $4.000000\ldots$ & $0$ \\
$+$\_real\_$D_4$                   & $1$ & $4.000000\ldots$ & $0$ \\
$x^4-2$ ($D_4$)                    & $1$ & $4.000000\ldots$ & $0$ \\
$x^4+1$ ($V_4$)                    & $1$ & $4.000000\ldots$ & $0$ \\
$x^4-x-1$ ($S_4$)                  & $1$ & $4.000000\ldots$ & $0$ \\
$\alpha_4=7$ sample                & $7$ & $2.45915261180505746086535722212\ldots$ & $0$ \\
\hline
\end{tabular}

The inferred constant
\[
c(4)\ :=\ \xi_0(b)\,\cdot\,a_4^{1/4}\ =\ 4
\]
is constant across all $8$ representatives to better than $10^{-50}$
(spread $0$).  \textbf{Claim Q1-B verdict: SUPPORTED}, $c(4)=4$.

\subsection*{General-$d$ formula}

By the same Newton-polygon derivation,
\[
\boxed{\ \xi_0(b)\ =\ \frac{d}{\beta_d^{1/d}}\quad\text{for all }d\ge 1,\ }
\]
in particular $c(d)=d$.  This recovers Proposition~3.3.A of
\textsc{channel-theory-v11} ($\xi_0=2/\sqrt{\beta_2}$ at $d=2$) and
predicts $\xi_0=3/\beta_3^{1/3}$ at $d=3$. The identity is an
\emph{operational} characteristic root of the irregular singularity,
not yet a Borel-plane radius identity (the latter remains open
per CC-PIPELINE-G).
